\usepackage{fontspec}
\usepackage{polyglossia}
\setdefaultlanguage{vietnamese}
\setmainfont{texgyretermes}[
    Extension      = .otf,
    UprightFont    = *-regular,
    BoldFont       = *-bold,
    ItalicFont     = *-italic,
    BoldItalicFont = *-bolditalic
]
\usepackage[a4paper,margin=2.5cm]{geometry}
\usepackage{setspace}
\onehalfspacing
\usepackage{indentfirst}
\setlength{\emergencystretch}{3em}
\usepackage{titlesec}
\usepackage{fancyhdr}
\usepackage{graphicx}
\usepackage{booktabs}
\usepackage{longtable}
\usepackage{array}
\usepackage{float}
\usepackage{amsmath,amssymb}
\usepackage{xcolor}
\usepackage{enumitem}
\usepackage{caption}
\usepackage{hyperref}
\hypersetup{plainpages=false, colorlinks=true, linkcolor=black, urlcolor=blue!60!black, citecolor=black}

\newcommand{\schoolname}{Trường Đại học Bách Khoa -- ĐHQG-HCM}
\newcommand{\facultyname}{Khoa Khoa học và Kỹ thuật Máy tính}
\newcommand{\coursename}{Nền tảng lập trình cho phân tích và trực quan dữ liệu}
\newcommand{\coursecode}{CO5177}
\newcommand{\semestername}{Học kỳ 261, năm học 2026--2027}
\newcommand{\lecturername}{TS. Lê Thành Sách}
\newcommand{\groupname}{G1}
\newcommand{\landingurl}{https://minhtho1503.github.io/data-analysis-visualization-g1/}
\newcommand{\repourl}{https://github.com/minhtho1503/data-analysis-visualization-g1}
\newcommand{\reportdate}{TP. Hồ Chí Minh, tháng 11 năm 2026}
\newcommand{\studentrows}{%
  HVTH: & Phan Thị Minh Thơ -- 2470112 \\
        & Nguyễn Thị Thu An -- 2470315 \\
}

\pagestyle{fancy}
\fancyhf{}
\lhead{\leftmark}
\rhead{\thepage}
\fancypagestyle{plain}{\fancyhf{}\rhead{\thepage}\renewcommand{\headrulewidth}{0pt}}

\titleformat{\chapter}[display]{\normalfont\huge\bfseries}{\chaptertitlename\ \thechapter}{20pt}{\Huge}
\titleformat{\section}{\normalfont\Large\bfseries}{\thesection}{1em}{}
\titleformat{\subsection}{\normalfont\large\bfseries}{\thesubsection}{1em}{}

\newcommand{\reporttype}{BÁO CÁO BÀI TẬP LỚN}
\newcommand{\reportsubtitle}{}

\newcommand{\figdir}{}
\newcommand{\reportfig}[3]{%
  \begin{figure}[H]\centering
    \IfFileExists{\figdir#1}%
      {\includegraphics[width=\linewidth]{\figdir#1}}%
      {\fbox{\parbox[c][5cm][c]{0.9\linewidth}{\centering\ttfamily chưa có hình: \detokenize{#1}}}}%
    \caption{#2}\label{#3}
  \end{figure}}
